% GENERATED FILE. Edit canonical lessons, then run npm run generate:books.
% canonical-source-hash: fnv1a64:02a1fee4fa7fd2a7
% canonical-lessons: ES-C381-lavar, ES-C381-romper, ES-C381-saltar, ES-C381-sonar

\chapter{Four Everyday Actions English Kept}
\label{ch:lavar-y-romper}
\hlchaptermodality{381}

\begin{chapteropening}
\textbf{By the end of this chapter:} \emph{I can say lavar, romper, saltar and soñar, and name four everyday actions.}

Complete ES-C381-sonar: name four everyday actions, and say which one you do most often.
\end{chapteropening}

\section[lavar]{lavar — a washing-place, and a downpour}

\label{lesson:ES-C381-lavar}



\begin{quote}
Say \emph{last}, then \emph{necessary}. (\emph{Último}; \emph{necesario}.)
\end{quote}

\subsection*{What to know first}
\begin{itemize}
\raggedright
  \item el agua — what you wash with.
  \item la cara — one thing you wash.
\end{itemize}

\begin{sounds}
\begin{itemize}
\raggedright
  \item \emph{la-VAR}, two beats with the weight on the second. The \emph{v} is the same sound as \emph{b}, and the final \emph{r} is a single tap. $\to$ pronunciation reference
\end{itemize}
\end{sounds}

\begin{cousinweb}
\textbf{lavar} — \textquotedblleft{}to wash.\textquotedblright{}

Latin \textbf{lavare} meant to wash, and that is the whole of the Spanish story. Said like that it is a circle: the word for washing comes from the word for washing. So the interesting part is not where \emph{lavar} came from. It is where else the same Latin verb went.

English has it three times over, and never looking like itself.

A \textbf{lavatory} is, literally, a washing-place. The room got its name from the basin, not the other way round — which is why the word sounds so formal for so ordinary a thing.

To \textbf{launder} is the same verb, arrived by a longer road: Latin \emph{lavandaria}, things to be washed, through French, losing a syllable on the way. The modern sense of laundering money is a metaphor somebody built on top of an already buried one.

And \textbf{lavish} is the surprise. Before it meant generous, it meant a downpour — a \emph{lavasse}, a sudden heavy rain. Spending lavishly is spending the way rain falls.

One caution, because it is the connection everyone reaches for first: \textbf{lava} and \textbf{lavender} are \emph{not} safe to add to this list. Both have disputed etymologies, and the tidy story that lava is \textquotedblleft{}what washes down the mountain\textquotedblright{} is one proposal among several rather than the answer.
\end{cousinweb}

\begin{grammarlens}[title={What you've built}]
You can now name the action: \emph{lavar}. And you can hear it inside an English word you already knew.
\end{grammarlens}

\subsection*{Your turn}
Say these aloud:

\begin{itemize}
\raggedright
  \item \emph{lavar}
  \item \emph{lavar la cara}
  \item \emph{el agua}
\end{itemize}

\begin{tcolorbox}[breakable,colback=teal!4,colframe=teal!35!black,title={Before you move on}]
Which English room is named for its basin? (\emph{Lavatory}.) And what did \emph{lavish} mean before it meant generous? (A downpour.)
\end{tcolorbox}
\section[romper]{romper — a road is something broken through a forest}

\label{lesson:ES-C381-romper}



\begin{quote}
Say \emph{necessary}. (\emph{Necesario}.) And \emph{to wash}? (\emph{Lavar}.)
\end{quote}

\subsection*{What to know first}
\begin{itemize}
\raggedright
  \item la puerta — something breakable.
  \item la ventana — another.
\end{itemize}

\begin{sounds}
\begin{itemize}
\raggedright
  \item \emph{rom-PER}, two beats with the weight on the second. The opening \emph{r} is the strong trilled one at the start of a word. $\to$ pronunciation reference
\end{itemize}
\end{sounds}

\begin{cousinweb}
\textbf{romper} — \textquotedblleft{}to break.\textquotedblright{}

Latin \textbf{rumpere}, to break or burst. Again the direct line is a circle, so take the branch instead — and this one is worth the walk.

The past participle of \emph{rumpere} was \textbf{rupta}, \textquotedblleft{}broken.\textquotedblright{} Roman Latin had the phrase \textbf{rupta via}: a \emph{broken way}. It meant a path forced through a forest by breaking the trees down, as opposed to a road that was built.

Drop the \emph{via}, keep the adjective, and \emph{rupta} becomes French \emph{route} and then English \textbf{route}. Every route you have ever taken is, etymologically, something somebody smashed a way through.

English also kept the verb straight, in \textbf{rupture}, \textbf{erupt} and \textbf{abrupt} — an abrupt ending is one that breaks off.

A false friend to refuse: \textbf{romp} has nothing to do with \emph{romper}. It is a much later English word of unclear origin, and the resemblance is an accident.
\end{cousinweb}

\begin{grammarlens}[title={What you've built}]
You can now name the action: \emph{romper}. And you know why a route is not a built road.
\end{grammarlens}

\subsection*{Your turn}
Say these aloud:

\begin{itemize}
\raggedright
  \item \emph{romper}
  \item \emph{romper la ventana}
  \item \emph{lavar la puerta}
\end{itemize}

\begin{tcolorbox}[breakable,colback=teal!4,colframe=teal!35!black,title={Before you move on}]
What was a \emph{rupta via}? (A way broken through a forest.) Is English \emph{romp} related? (No — an accident.)
\end{tcolorbox}
\section[saltar]{saltar — a jump over the top}

\label{lesson:ES-C381-saltar}



\begin{quote}
Say \emph{to break}. (\emph{Romper}.) And \emph{to wash}? (\emph{Lavar}.)
\end{quote}

\subsection*{What to know first}
\begin{itemize}
\raggedright
  \item el pie — what you jump with.
  \item arriba — where you jump to.
\end{itemize}

\begin{sounds}
\begin{itemize}
\raggedright
  \item \emph{sal-TAR}, two beats with the weight on the second. The \emph{l} is clear, never swallowed the way English swallows it in \emph{salt}. $\to$ pronunciation reference
\end{itemize}
\end{sounds}

\begin{cousinweb}
\textbf{saltar} — \textquotedblleft{}to jump.\textquotedblright{}

Latin \textbf{saltare} meant to leap, and also to dance — the two were one idea to a Roman. The Spanish is the Latin, so once more the branch is the interesting part.

The best one is \textbf{somersault}, which does not look like it contains a Latin verb at all. It is \emph{supra saltus} — a jump \emph{over the top}. The first half is Latin \emph{supra}, above, the same \emph{supra} in \emph{supranational}; the second half is this verb. A somersault is a going-over-jump, and the word has been carrying that description, unread, for centuries.

Two more: \textbf{assault} is \emph{ad-} plus this verb, leaping \emph{at} somebody. And a \textbf{sally}, when troops rush out of a fort, is the same leap from the French.

The false friend is right there in the pronunciation note: \textbf{salt} is not related. English \emph{salt} is a native Germanic word, cousin to Latin \emph{sal}, and has nothing to do with jumping. \emph{Salt} and \emph{somersault} only rhyme.
\end{cousinweb}

\begin{grammarlens}[title={What you've built}]
You can now name the action: \emph{saltar}. And you can take a somersault apart into two Latin words.
\end{grammarlens}

\subsection*{Your turn}
Say these aloud:

\begin{itemize}
\raggedright
  \item \emph{saltar}
  \item \emph{saltar arriba}
  \item \emph{romper}
\end{itemize}

\begin{tcolorbox}[breakable,colback=teal!4,colframe=teal!35!black,title={Before you move on}]
What are the two halves of \emph{somersault}? (\emph{Supra} — above; \emph{saltus} — a jump.) Is \emph{salt} related? (No.)
\end{tcolorbox}
\section[soñar]{soñar — the tilde used to be an m}

\label{lesson:ES-C381-sonar}



\begin{quote}
Say \emph{to jump}. (\emph{Saltar}.) And \emph{to break}? (\emph{Romper}.)
\end{quote}

\subsection*{What to know first}
\begin{itemize}
\raggedright
  \item la cama — where you do it.
  \item la noche — when.
\end{itemize}

\begin{sounds}
\begin{itemize}
\raggedright
  \item \emph{so-ÑAR}, two beats with the weight on the second. The \emph{ñ} is one sound, the \emph{ny} of \emph{canyon}, not an \emph{n} followed by a \emph{y}. $\to$ pronunciation reference
\end{itemize}
\end{sounds}

\begin{cousinweb}
\textbf{soñar} — \textquotedblleft{}to dream.\textquotedblright{}

Latin \textbf{somnus} was sleep; \textbf{somniare} was to dream. Spanish took the verb.

Look at what happened to the \textbf{m}. \emph{Somniare} has \emph{mn} in the middle, and Spanish would not keep that cluster — the two nasals collapsed into a single new sound, and the letter invented to write it was \textbf{ñ}. The tilde itself began as a small \emph{n} written above the line by scribes saving parchment. So the mark on top of \emph{soñar} is, quite literally, the second nasal of the Latin word, moved upstairs.

English kept \emph{somnus} rather than the verb. \textbf{Insomnia} is \emph{in-} plus \emph{somnus}: not-sleep. \textbf{Somnolent} means sleepy, and a \textbf{somnambulist} is a sleep-walker — \emph{somnus} plus \emph{ambulare}, to walk.

A false friend, and an easy one to fall into: \textbf{sonar} without the tilde is a different Spanish verb, meaning to sound, and English \emph{sonar} is not Spanish at all — it is a modern acronym. The tilde is doing real work here.
\end{cousinweb}

\begin{grammarlens}[title={What you've built}]
You can now name four everyday actions: \emph{lavar}, \emph{romper}, \emph{saltar}, \emph{soñar}. Say which one you do most often.
\end{grammarlens}

\subsection*{Your turn}
Say these aloud:

\begin{itemize}
\raggedright
  \item \emph{soñar}
  \item \emph{soñar en la cama}
  \item \emph{lavar, romper, saltar, soñar}
\end{itemize}

\begin{tcolorbox}[breakable,colback=teal!4,colframe=teal!35!black,title={Before you move on}]
What does \emph{insomnia} take apart into? (\emph{In-} not, \emph{somnus} sleep.) And what did the tilde on \emph{ñ} used to be? (A small \emph{n} written above the line.)
\end{tcolorbox}
